\renewcommand{\tema}{Examen Diagnóstico - Trabajo y Leyes de la Conservación}

\twocolumn
\section{Instrucciones}

Dispones de \textbf{15 minutos} para resolver 10 preguntas.



\section{PREGUNTAS DIAGNÓSTICO}

\begin{preguntabox}
\subsection{Pregunta 1}
¿Cuál de las siguientes situaciones representa trabajo mecánico \textbf{nulo}?

\begin{enumerate}
    \item[A)] Un cargador levanta una caja del suelo hasta su hombro.
    \item[B)] Un estudiante empuja una pared con toda su fuerza sin moverla.
    \item[C)] Una pelota cae libremente desde cierta altura.
    \item[D)] Un auto frena y reduce su velocidad a la mitad.
\end{enumerate}
\end{preguntabox}

\begin{preguntabox}
\subsection{Pregunta 2}
La energía cinética de un objeto depende de su masa y de su velocidad. Si la velocidad de un cuerpo se \textbf{duplica} y su masa permanece constante, ¿qué ocurre con su energía cinética?

\begin{enumerate}
    \item[A)] Se duplica.
    \item[B)] Se triplica.
    \item[C)] Se cuadruplica.
    \item[D)] Permanece igual.
\end{enumerate}
\end{preguntabox}

\begin{preguntabox}
\subsection{Pregunta 3}
Una persona jala una caja con una fuerza de 100 N formando un ángulo de 60$^\circ$ con la horizontal, desplazándola 5 m. ¿Cuánto trabajo realizó?\\
Considera $\cos 60$$^\circ$ $= 0.5$

\begin{enumerate}
    \item[A)] 25 J
    \item[B)] 250 J
    \item[C)] 500 J
    \item[D)] 866 J
\end{enumerate}
\end{preguntabox}
\newpage
\begin{preguntabox}
\subsection{Pregunta 4}
Un automóvil se mueve a velocidad constante de 20 m/s venciendo una fuerza de rozamiento de 400 N. ¿Cuál es la potencia que desarrolla el motor?

\begin{enumerate}
    \item[A)] 20 W
    \item[B)] 420 W
    \item[C)] 4\,000 W
    \item[D)] 8\,000 W
\end{enumerate}
\end{preguntabox}


\begin{preguntabox}
\subsection{Pregunta 5}
Un motor de 500 W impulsa un vehículo aplicando una fuerza constante de 250 N. ¿Con qué velocidad se mueve el vehículo?

\begin{enumerate}
    \item[A)] 0.5 m/s
    \item[B)] 2 m/s
    \item[C)] 250 m/s
    \item[D)] 125\,000 m/s
\end{enumerate}
\end{preguntabox}

\begin{preguntabox}
\subsection{Pregunta 6}
Un objeto en movimiento tiene una energía cinética de 200 J y se desplaza a 10 m/s. ¿Cuál es su masa?

\begin{enumerate}
    \item[A)] 2 kg
    \item[B)] 4 kg
    \item[C)] 20 kg
    \item[D)] 40 kg
\end{enumerate}
\end{preguntabox}
\newpage
\begin{preguntabox}
\subsection{Pregunta 7}
Se coloca un mismo bloque en la cima de tres rampas de igual altura $h$ pero diferente forma, como se muestra. ¿Cuál de las siguientes afirmaciones es correcta respecto a la energía potencial gravitatoria del bloque en cada caso?

\begin{center}
\shorthandoff{>}\shorthandoff{<}
\begin{tikzpicture}[scale=0.72]
  % --- Rampa I: semicírculo ---
  \fill[black!80] (0,0) arc[start angle=180, end angle=0, radius=1cm] -- cycle;
  \fill[gray!60] (0.6, 1.0) rectangle (1.4, 1.35);
  \draw[thick] (-1.5,0) -- (2.5,0);
  % cota h
  \draw[<->,>=Stealth, thin] (-1.2,0) -- (-1.2,1.0) node[midway, left] {\small $h$};
  \node at (0.5,-0.4) {\small I};

  % --- Rampa II: pedestal recto ---
  \fill[black!80] (3.5,0) rectangle (4.5, 1.0);
  \fill[gray!60] (3.7, 1.0) rectangle (4.3, 1.35);
  \draw[thick] (2.8,0) -- (5.2,0);
  \node at (4.0,-0.4) {\small II};

  % --- Rampa III: triángulo ---
  \fill[black!80] (6.0,0) -- (7.0,1.0) -- (8.0,0) -- cycle;
  \fill[gray!60] (6.8, 1.0) rectangle (7.2, 1.35);
  \draw[thick] (5.5,0) -- (8.5,0);
  \node at (7.0,-0.4) {\small III};
\end{tikzpicture}
\shorthandon{>}\shorthandon{<}
\end{center}

\begin{enumerate}
    \item[A)] El bloque en la rampa III tiene mayor energía potencial porque la pendiente es más pronunciada.
    \item[B)] El bloque en la rampa I tiene menor energía potencial porque la superficie es curva.
    \item[C)] Los tres bloques tienen la misma energía potencial gravitatoria.
    \item[D)] La energía potencial depende de la forma de la rampa, por lo que los tres casos son distintos.
\end{enumerate}
\end{preguntabox}

\begin{preguntabox}
\subsection{Pregunta 8}
Una roca de 5 kg se suelta desde el reposo a 20 m de altura. Despreciando la fricción del aire, ¿con qué velocidad llega al suelo?\\
Considera $g = 10$ m/s$^2$\\
\textit{(Resuelve usando conservación de energía, no cinemática.)}

\begin{enumerate}
    \item[A)] 10 m/s
    \item[B)] 20 m/s
    \item[C)] 200 m/s
    \item[D)] 400 m/s
\end{enumerate}
\end{preguntabox}

\begin{preguntabox}
\subsection{Pregunta 9}
Dos patinadores en reposo sobre hielo (sin fricción) se empujan mutuamente. El patinador A tiene 70 kg y el patinador B tiene 35 kg. Si B sale con velocidad de 4 m/s, ¿con qué velocidad sale A?

\begin{enumerate}
    \item[A)] 2 m/s en la misma dirección que B.
    \item[B)] 4 m/s en dirección opuesta a B.
    \item[C)] 2 m/s en dirección opuesta a B.
    \item[D)] 8 m/s en dirección opuesta a B.
\end{enumerate}
\end{preguntabox}

\begin{preguntabox}
\subsection{Pregunta 10}
Un bloque se desliza por un plano con fricción y llega al fondo con menos energía cinética que la energía potencial que tenía en la cima. ¿Qué ocurrió con la energía ``perdida''?

\begin{enumerate}
    \item[A)] Se destruyó; la energía no siempre se conserva.
    \item[B)] Se transformó en energía térmica por la fricción.
    \item[C)] Quedó almacenada como energía potencial en el plano.
    \item[D)] Se convirtió en ímpetu del sistema.
\end{enumerate}
\end{preguntabox}

\vspace{1cm}

\begin{tcolorbox}[colback=amarillonota, colframe=orange, boxrule=0pt, leftrule=3pt]
\textbf{Detén el cronómetro.} Pasa a la siguiente sección para revisar tus respuestas.
\end{tcolorbox}

\newpage
\onecolumn
\section{RESPUESTAS Y RETROALIMENTACIÓN}

\begin{respuestabox}
\subsection{Pregunta 1}
\textcolor{verdecorrecto}{\textbf{Respuesta correcta: B) Un estudiante empuja una pared sin moverla.}}

\textbf{Justificación:}\\
El trabajo mecánico se define como:
$$W = F \cdot d \cdot \cos\theta$$
Para que exista trabajo, debe haber desplazamiento ($d \neq 0$). Si la pared no se mueve, $d = 0$, por lo tanto $W = 0$, sin importar cuánta fuerza se aplique.

\textbf{Respuestas incorrectas:}
\begin{itemize}
    \item \textbf{A) El cargador levanta la caja:} Hay fuerza y desplazamiento en la misma dirección $\Rightarrow W > 0$. Error: confundir esfuerzo físico con trabajo nulo.
    \item \textbf{C) La pelota cae libremente:} La gravedad actúa en el mismo sentido del movimiento $\Rightarrow W > 0$. Error: pensar que lo ``natural'' implica trabajo nulo.
    \item \textbf{D) El auto frena:} La fricción actúa en sentido contrario al movimiento $\Rightarrow W < 0$ (trabajo negativo, no nulo).
\end{itemize}
\end{respuestabox}

\begin{respuestabox}
\subsection{Pregunta 2}
\textcolor{verdecorrecto}{\textbf{Respuesta correcta: C) Se cuadruplica.}}

\textbf{Justificación:}\\
La energía cinética es:
$$E_c = \frac{1}{2}mv^2$$
La velocidad aparece al \textbf{cuadrado}. Si $v' = 2v$, entonces:
$$E_c' = \frac{1}{2}m(2v)^2 = \frac{1}{2}m \cdot 4v^2 = 4 \cdot E_c$$

\textbf{Respuestas incorrectas:}
\begin{itemize}
    \item \textbf{A) Se duplica:} Error de ignorar el exponente. Se duplicaría si la relación fuera $E_c \propto v$, pero es $E_c \propto v^2$.
    \item \textbf{B) Se triplica:} No corresponde a ninguna relación matemática de la fórmula.
    \item \textbf{D) Permanece igual:} Implicaría que la energía es independiente de la velocidad, lo cual contradice la definición.
\end{itemize}
\end{respuestabox}

\begin{respuestabox}
\subsection{Pregunta 3}
\textcolor{verdecorrecto}{\textbf{Respuesta correcta: B) 250 J}}

\textbf{Justificación:}\\
Cuando la fuerza forma un ángulo con el desplazamiento, solo la componente horizontal realiza trabajo:
$$W = F \cdot d \cdot \cos\theta = 100\,\text{N} \times 5\,\text{m} \times \cos 60^\circ = 500 \times 0.5 = 250\,\text{J}$$

\textbf{Respuestas incorrectas:}
\begin{itemize}
    \item \textbf{A) 25 J:} Error de calcular $\frac{F \cdot \cos\theta}{d} = \frac{50}{2}$. Confusión de operaciones.
    \item \textbf{C) 500 J:} Error de ignorar el ángulo ($\cos 60$$^\circ$). Usa $W = F \cdot d$ sin el factor coseno.
    \item \textbf{D) 866 J:} Usa $\sin 60$$^\circ$ $\approx 0.866$ en lugar de $\cos 60$$^\circ$. Error de componente: el seno da la componente perpendicular al movimiento, no la que realiza trabajo.
\end{itemize}
\end{respuestabox}

\newpage

\begin{respuestabox}
\subsection{Pregunta 4}
\textcolor{verdecorrecto}{\textbf{Respuesta correcta: D) 8\,000 W}}

\textbf{Justificación:}\\
A velocidad constante, la fuerza del motor es igual a la fuerza de rozamiento. Se aplica la relación:
$$P = F \cdot v = 400\,\text{N} \times 20\,\text{m/s} = 8\,000\,\text{W}$$

\textbf{Respuestas incorrectas:}
\begin{itemize}
    \item \textbf{A) 20 W:} Resulta de dividir $v \div F = 20 \div 400$. Error de invertir la operación.
    \item \textbf{B) 420 W:} Resulta de sumar $F + v = 400 + 20$. Error de confundir la operación.
    \item \textbf{C) 4\,000 W:} Resulta de $\frac{F \times v}{2} = \frac{8000}{2}$. Posible confusión con la fórmula de energía cinética que lleva el factor $\frac{1}{2}$.
\end{itemize}
\end{respuestabox}

\begin{respuestabox}
\subsection{Pregunta 5}
\textcolor{verdecorrecto}{\textbf{Respuesta correcta: B) 2 m/s}}

\textbf{Justificación:}\\
Se despeja $v$ de la relación $P = F \cdot v$:
$$v = \frac{P}{F} = \frac{500\,\text{W}}{250\,\text{N}} = 2\,\text{m/s}$$

\textbf{Respuestas incorrectas:}
\begin{itemize}
    \item \textbf{A) 0.5 m/s:} Resulta de invertir el despeje: $v = F/P = 250/500$. Error de dividir al revés.
    \item \textbf{C) 250 m/s:} Usa solo la fuerza como si fuera la velocidad. Confusión de variables.
    \item \textbf{D) 125\,000 m/s:} Resulta de multiplicar $P \times F = 500 \times 250$. Error de aplicar producto en lugar de cociente.
\end{itemize}
\end{respuestabox}

\begin{respuestabox}
\subsection{Pregunta 6}
\textcolor{verdecorrecto}{\textbf{Respuesta correcta: B) 4 kg}}

\textbf{Justificación:}\\
Se despeja $m$ de la fórmula de energía cinética:
$$E_c = \frac{1}{2}mv^2 \quad \Rightarrow \quad m = \frac{2E_c}{v^2} = \frac{2(200)}{(10)^2} = \frac{400}{100} = 4\,\text{kg}$$

\textbf{Respuestas incorrectas:}
\begin{itemize}
    \item \textbf{A) 2 kg:} Error de no aplicar el factor 2 al despejar: $m = E_c / v^2 = 200/100 = 2$. Olvida el $2E_c$.
    \item \textbf{C) 20 kg:} Error de no elevar la velocidad al cuadrado: $m = 2E_c / v = 400/10 = 40$. Calcula mal $v^2$.
    \item \textbf{D) 40 kg:} Error de no elevar al cuadrado \textit{y} no aplicar el factor 2: $m = E_c/v = 200/10 = 20$, o bien $2E_c/v = 400/10 = 40$.
\end{itemize}
\end{respuestabox}

\newpage

\begin{respuestabox}
\subsection{Pregunta 7}
\textcolor{verdecorrecto}{\textbf{Respuesta correcta: C) Los tres bloques tienen la misma energía potencial gravitatoria.}}

\textbf{Justificación:}\\
La energía potencial gravitatoria solo depende de la masa, la gravedad y la \textbf{altura}:
$$E_p = mgh$$
Como los tres bloques tienen la misma masa $m$ y están a la misma altura $h$, su energía potencial es idéntica en los tres casos. La forma de la rampa (curva, recta o inclinada) no interviene en absoluto en esta fórmula.

\textbf{Respuestas incorrectas:}
\begin{itemize}
    \item \textbf{A) Rampa III tiene mayor $E_p$ por la pendiente:} Error conceptual grave. La pendiente afecta la fuerza necesaria para subir, pero no la energía potencial almacenada. Solo importa la altura final.
    \item \textbf{B) Rampa I tiene menor $E_p$ por la superficie curva:} Error similar. La forma de la trayectoria es irrelevante para $E_p$; lo único que cuenta es la diferencia de altura respecto al nivel de referencia.
    \item \textbf{D) La forma de la rampa determina la $E_p$:} Confunde la \textit{trayectoria} con la \textit{posición final}. La energía potencial es una función de estado: solo depende de dónde está el objeto, no de cómo llegó ahí.
\end{itemize}
\end{respuestabox}

\begin{respuestabox}
\subsection{Pregunta 8}
\textcolor{verdecorrecto}{\textbf{Respuesta correcta: B) 20 m/s}}

\textbf{Justificación:}\\
Usando conservación de energía mecánica: toda la energía potencial se convierte en cinética:
$$E_p = E_c \quad \Rightarrow \quad mgh = \frac{1}{2}mv^2$$
La masa se cancela en ambos lados:
$$gh = \frac{1}{2}v^2 \quad \Rightarrow \quad v = \sqrt{2gh} = \sqrt{2(10)(20)} = \sqrt{400} = 20\,\text{m/s}$$
Nota: la cinemática daría el mismo resultado ($v^2 = v_0^2 + 2ad$), pero la conservación de energía es más directa y general.

\textbf{Respuestas incorrectas:}
\begin{itemize}
    \item \textbf{A) 10 m/s:} Error de omitir el factor 2: $v = \sqrt{gh} = \sqrt{200} \approx 14$ m/s (redondeado incorrectamente a 10).
    \item \textbf{C) 200 m/s:} Error de no extraer la raíz cuadrada al final: $2gh = 400$, pero $v = \sqrt{400}$, no $v = 400$.
    \item \textbf{D) 400 m/s:} Confunde $v^2 = 400$ con $v = 400$. El valor calculado es la velocidad \textit{al cuadrado}.
\end{itemize}
\end{respuestabox}
\newpage
\begin{respuestabox}
\subsection{Pregunta 9}
\textcolor{verdecorrecto}{\textbf{Respuesta correcta: C) 2 m/s en dirección opuesta a B.}}

\textbf{Justificación:}\\
El sistema parte del reposo, por lo que el ímpetu total inicial es cero. Por conservación del momento lineal:
$$m_A v_A + m_B v_B = 0$$
$$(70)\,v_A + (35)(4) = 0 \quad \Rightarrow \quad v_A = \frac{-140}{70} = -2\,\text{m/s}$$
El signo negativo indica dirección \textbf{opuesta} a B.

\textbf{Respuestas incorrectas:}
\begin{itemize}
    \item \textbf{A) 2 m/s en la misma dirección que B:} La magnitud es correcta, pero la dirección es imposible. Si ambos salieran en la misma dirección, el ímpetu total no sería cero, lo cual violaría la ley de conservación.
    \item \textbf{B) 4 m/s en dirección opuesta:} Dirección correcta, pero magnitud errónea. Solo sería válido si $m_A = m_B$; como A es el doble de pesado, su velocidad es la mitad.
    \item \textbf{D) 8 m/s en dirección opuesta:} Error de multiplicar las masas en lugar de dividirlas: $v_A = v_B \times (m_B / m_A) \times 2$.
\end{itemize}
\end{respuestabox}


\begin{respuestabox}
\subsection{Pregunta 10}
\textcolor{verdecorrecto}{\textbf{Respuesta correcta: B) Se transformó en energía térmica por la fricción.}}

\textbf{Justificación:}\\
La \textbf{Ley de Conservación de la Energía} establece que la energía total siempre se conserva. En presencia de fricción, parte de la energía mecánica se convierte en calor:
$$E_{m,i} = E_{m,f} + Q_{\text{fricción}}$$
donde $Q_{\text{fricción}} = f \cdot d$ es la energía disipada como calor.

\textbf{Respuestas incorrectas:}
\begin{itemize}
    \item \textbf{A) Se destruyó:} Contradice la Ley de Conservación de la Energía. La energía \textit{nunca} se destruye ni se crea.
    \item \textbf{C) Quedó como energía potencial:} La energía potencial depende de la altura, no de la fricción del trayecto.
    \item \textbf{D) Se convirtió en ímpetu:} El ímpetu (momento lineal) no es una forma de energía; son magnitudes distintas con unidades diferentes.
\end{itemize}
\end{respuestabox}